\section{Implemented source charts and reconstruction checks}
\label{sec:source-chart-implementation}

The implementation of Section~\ref{sec:coordinates} lives in
\path{Kernel/SourceCoordinates.wl}. The two chart families
are distinguished by the function used to reconstruct the source:
\wl{"SourceExp"} reconstructs an exponential, while
\wl{"SourceLog"} reconstructs a logarithm. The underlying inverse always
has observable power $1$. Powers of the original inverse are formed only
after that chart variable has been computed.

\paragraph{Pure logarithmic source dependence.}
First choose the positive local source coordinate $u$ from the endpoint
and direction. Substitute $u=e^{-h}$, so $h\to+\infty$. A pure
logarithmic expression becomes a function $\Phi(h)$ with no residual
exponential dependence on $h$, and the existing inverse dispatcher solves
$\Phi(h)=y$. For a finite endpoint $x_0$ and side $s\in\{-1,1\}$,
\[
 x=x_0+s e^{-h};
\]
at the source infinity with sign $s$, the reconstruction is $x=s e^h$.
Thus the noninteger real powers $(-\log u)^q$, polynomial and rational
combinations of $\log u$, and both finite source directions can use the
same chart whenever the transformed inverse belongs to a supported scale.

For example, $F(u)=(\log u)^2+\log u$ becomes $h^2-h$, and its small
positive inverse begins
\[
 u=e^{-\sqrt y-1/2}
 \left(1-\frac1{8\sqrt y}+\OO(y^{-1})\right).
\]
The exponential factor multiplies the absolute remainder. A finite-order
approximation to $h$ is usable here only when its absolute error tends to
zero. A finite relative Lambert expansion for a growing $h$ need not have
this property; increasing its logarithmic depth alone does not resolve
that obstruction.

An exact transformed inverse may instead be kept as an exact carrier.
Repeated exact charts therefore give, for instance, the exact inverse
$u=\exp(-e^y)$ of $\log(-\log u)$. Recursion is bounded to avoid an
uncontrolled recognizer loop. Exactness is recorded separately: it can
come from an exact inner series, an exact inner inverse, or a bounded
symbolic check that a small candidate composes exactly to the target.
The asymptotic branch of the inner inverse fixes which exact solution is
being reconstructed.

\paragraph{Finite exponential families.}
At $x\to s\infty$, the implementation takes $u=e^{-s x}$, so
$x=-s\log u$. Finite sums of real exponential rates become finite
generalized-power germs in $u$. Equal rates and cancellations are handled
by the ordinary canonicalization before a leading term is selected.
Rates need not be commensurable. The recognizer requires a nonzero leading
algebraic exponent after removing a constant target offset; additional
logarithmic factors are permitted when the resulting finite power--log
germ is supported by the inner engine.

For example,
\[
 e^{-x}+e^{-2x}=y
 \quad\Longrightarrow\quad
 x=-\log y+y-\frac32y^2+\frac{10}3y^3+\OO(y^4),
 \qquad y\to0^+.
\]
The inverse of $e^x+e^{-x}$ at $+\infty$ begins
$\log y-y^{-2}$; selecting the source endpoint $-\infty$ gives its
negative. The chart preserves this distinction rather than choosing a
numerical root without an endpoint condition.

\paragraph{Original observables and term counts.}
At a finite source endpoint, the power option represents
$(x-x_0)^r$ when $r\ne1$, and the inverse itself when $r=1$.
At infinity it represents $x^r$. A negative source side therefore requires
integer $r$ for these real-valued observable conventions. Exponential
reconstruction multiplies $h$ by $\pm r$ before exponentiating; a finite
source offset is added separately when $r=1$.

For a non-polynomial power of a logarithmic reconstruction, the finite
approximation can be retained inside a specified observable carrier.
If $X\sim c|\log w|$ and $x-X=\OO(w^P\M^D)$ with $P>0$, the mean
value theorem gives
\[
 \frac{x^r-X^r}{X^r}
 =\OO\left(\frac{x-X}{X}\right)
 =\OO(w^P\M^D).
\]
This deliberately conservative relative bound makes the retained carrier
$X^r$ usable without claiming a finite polynomial--log expansion of its
fractional power. An explicit cutoff is required for that representation;
a unit-block term count would not count the corrections retained inside
the carrier. In the ordinary reconstructed unit, the term-goal option
counts the actual nonzero reconstructed blocks.

\paragraph{Residuals and numerical comparisons.}
The object records both directions of the chart and an expression that
recovers the chart variable from the requested observable. Formal
residual checking applies that expression to the \emph{returned finite
observable}. It does not reuse the unmodified inner inverse blocks after
an exponential or logarithmic reconstruction has introduced a new
truncation. The checked residual is
\[
 \frac{\Phi(h_{\mathrm{returned}})-y}{y-y_0},
\]
with denominator $y$ at an infinite target. If reconstruction requires a
coefficient algebra outside the supported jet scale, the exact residual
expression is still returned, accompanied by an explicit indication that
no jet-order assertion was made.

Numerical comparison solves $\Phi(h)=y$ at high precision, reconstructs
the source, checks its original side, and compares the requested observable
with the returned approximation. This avoids unnecessarily evaluating a
very large or very small original exponential during root finding. It is
numerical evidence, not an interval certificate. Refinement replays the
original inverse problem and reconstruction, preserving inverse provenance.
An externally supplied forward remainder must first be transported through
the source chart; such a declaration is rejected until that transport has
been specified.
